\section{Extension Experiment: VQA-RAD Medical VQA Hallucination Detection}
\label{sec:vqarad}

Medical visual question answering is a representative application
of MLLMs in high-stakes domains.
Compared with general VQA and mathematical reasoning,
medical VQA imposes higher demands on both visual perception
and domain knowledge,
requiring models to correctly identify anatomical structures
and pathological findings in X-ray, CT, and MRI images,
understand the clinical semantics of questions,
and provide diagnosis-related answers grounded in visual evidence.
This section evaluates the four models
on VQA-RAD \citep{lauDatasetClinicallyGenerated2018},
a standard radiology VQA benchmark.

\subsection{Results}

\Cref{tab:vqarad_results} reports the overall evaluation results
and hallucination rates stratified by answer type.

\begin{table}[htbp]
  \centering
  \caption{VQA-RAD evaluation results. Best values are highlighted in pink.}
  \label{tab:vqarad_results}
  \small
  \begin{tabular}{l cccc}
    \toprule
    \multicolumn{5}{c}{\textbf{Overall Results}} \\
    \midrule
    Metric
      & GPT-5.4-mini & Gemini 2.5 Flash & Qwen3.5-35B & Qwen3-VL-235B \\
    \midrule
    Hallucination Rate & \cellcolor{pink!25}33.26\% & 66.30\% & 39.25\% & 58.31\% \\
    Average Score      & \cellcolor{pink!25}4.52    & 2.68    & 4.20    & 3.57 \\
    $n_{\text{hallucinated}}$ & \cellcolor{pink!25}150 & 299 & 177 & 263 \\
    \quad Faithfulness & 145 (96.7\%) & 296 (99.0\%) & 171 (96.6\%) & 257 (97.7\%) \\
    \quad Factuality   & 1 (0.7\%)    & 2 (0.7\%)    & 4 (2.3\%)    & 5 (1.9\%) \\
    \quad Logical      & 4 (2.7\%)    & 1 (0.3\%)    & 2 (1.1\%)    & 1 (0.4\%) \\
    None               & 301 & 152 & 274 & 188 \\
    \midrule
    \multicolumn{5}{c}{\textbf{Stratified by Answer Type}} \\
    \midrule
    Answer Type ($N$)
      & GPT-5.4-mini & Gemini 2.5 Flash & Qwen3.5-35B & Qwen3-VL-235B \\
    \midrule
    Closed (yes/no, 251)   & \cellcolor{pink!25}28.69\% & 61.75\% & 31.87\% & 50.20\% \\
    Open (free-text, 200)  & \cellcolor{pink!25}39.00\% & 72.00\% & 48.50\% & 68.50\% \\
    \bottomrule
  \end{tabular}
\end{table}

%\begin{figure}[htbp]
%  \centering
%  \includegraphics[width=0.85\textwidth]{fig6_vqarad_pie.pdf}
%  \caption{Hallucination type composition for VQA-RAD.
%    Faithfulness dominates ($>$96\%) across all models.}
%  \label{fig:vqarad_pie}
%\end{figure}

\subsection{Analysis}

The VQA-RAD results reveal a performance landscape
that differs markedly from both MathVista and POPE.

\textbf{GPT-5.4-mini: strongly bimodal performance.}
GPT-5.4-mini leads overall
with the lowest HR (33.26\%) and highest average score (4.52)
across both closed and open questions.
Its score distribution, however, is strongly bimodal:
142 samples receive score 2
(informative but visually incorrect)
while 193 receive score 6 (perfect),
with intermediate scores comparatively rare.
Rather than being uniformly strong,
the model alternates between confident-and-correct
and confident-and-wrong,
consistent with the conservative response strategy
observed on POPE.

\textbf{Model ranking reversal.}
Qwen3-VL-235B (HR = 58.31\%), which leads on POPE,
and Qwen3.5-35B-A3B (HR = 39.25\%), which leads on MathVista,
both regress in the medical domain.
Conversely, GPT-5.4-mini,
which trails on MathVista and ranks among the weaker on POPE,
performs best here.
This reversal indicates that medical VQA
depends on a different axis of capability:
not general visual perception (probed by POPE)
nor chart-style reasoning (probed by MathVista),
but \emph{radiology-specific} visual perception:
the ability to localize anatomical structures
and pathological findings on X-ray, CT, and MRI inputs.
General-purpose MLLMs evidently have uneven coverage
of radiology imagery in their training data.

\textbf{Faithfulness dominates the error profile.}
Faithfulness hallucinations account for over 96\% of all hallucinations,
in contrast with MathVista
where all three types are well represented.
Almost all hallucinations stem from visual misreadings
rather than from logical or factual errors.
The core challenge of medical VQA
lies in precise identification
of anatomical structures and pathological findings:
whether the aorta is widened on a chest X-ray
(\Cref{sec:vqarad_case1}),
or which hemisphere of a brain MRI contains a lesion
(\Cref{sec:vqarad_case2}).
When the model misjudges the anatomical content of the image,
all downstream answers based on the incorrect premise
are classified as faithfulness hallucinations by the judge.

\subsection{Failure Case Analysis}

\subsubsection{Case 1: Conservative Refusal Leading to Aortic Aneurysm
  Missed Diagnosis (Closed Question, ID=0)}
\label{sec:vqarad_case1}

\Cref{tab:vqarad_case1_answers} and \Cref{fig:vqarad_case1}
present this case.
The question asks whether there is evidence
of an aortic aneurysm,
with the expected answer being Yes.

\begin{figure}[htbp]
  \centering
  \includegraphics[width=0.30\textwidth]{vqarad_fc_0.jpg}
  \caption{VQA-RAD Case 1: chest X-ray.}
  \label{fig:vqarad_case1}
\end{figure}

\begin{table}[htbp]
  \centering
  \caption{VQA-RAD Case 1: model responses.
    Expected answer: Yes (aortic aneurysm present).}
  \label{tab:vqarad_case1_answers}
  \small
  \begin{tabular}{p{2.8cm}p{8cm}cc}
    \toprule
    Model & Response summary & Score & Type \\
    \midrule
    GPT-5.4-mini     & No convincing evidence; mild aortic widening      & 2 & faithfulness \\
    Gemini 2.5 Flash & Limited sensitivity on chest X-ray, cannot confirm & 2 & faithfulness \\
    Qwen3.5-35B      & Evidence suggests descending aortic aneurysm      & 6 & none \\
    Qwen3-VL-235B    & No clear radiographic evidence                    & 2 & faithfulness \\
    \bottomrule
  \end{tabular}
\end{table}

\textbf{Three of four converge on under-claim.}
Three of the four models (GPT-5.4-mini, Gemini 2.5 Flash, Qwen3-VL-235B)
return some variant of ``no convincing evidence,''
even though all three describe the relevant finding (mild aortic
widening or prominence in the mediastinum) and simply stop short
of calling it a diagnosis.
Perception succeeds; inference quietly abstains.
Only Qwen3.5-35B-A3B converts the same observation
into the correct positive diagnosis and receives the highest score (6),
confirming that the radiographic signal
is in fact recoverable from the image.

\textbf{Conservative strategy generalizes from POPE.}
On POPE, GPT-5.4-mini shows a conservative response profile
(\Cref{sec:pope_cross_model_comparison}),
with the lowest Yes Ratio (0.414 on Random)
and highest Precision (0.9831 on Random)
at the cost of recall.
This profile extends to three of the four models
in the radiological setting.
A plausible driver is conservative alignment
on medically sensitive prompts:
the rehearsed safe phrasing
``no clear radiographic evidence; clinical correlation needed''
is precisely the kind of clinically defensible non-answer
that medical instruction-tuning data would reinforce.
Regardless of the exact training cause,
the empirical effect is a panel-wide under-claim bias
on closed-ended screening questions.

\textbf{Asymmetric costs in clinical screening.}
The clinical cost of this profile is sharper than its POPE counterpart.
Defaulting to No on POPE produces benign omissions
(\Cref{sec:pope_cross_model_comparison});
on a radiological screening question,
the same behavior becomes a missed diagnosis,
since the false negative blocks the patient
from entering the downstream diagnostic pathway,
which is the very purpose of a screening question.
A conservative response strategy therefore does not transfer
cleanly between general VQA and clinical screening:
in the latter setting,
the under-claim itself becomes the dangerous outcome,
rather than a safe default.

\subsubsection{Case 2: Medical Knowledge Overriding Visual Evidence
  (Open Question, ID=102)}
\label{sec:vqarad_case2}

\Cref{tab:vqarad_case2_answers} and \Cref{fig:vqarad_case2}
present this case.
The question asks on which side of the brain a lesion is located,
with the ground-truth answer ``left.''
On the axial MRI shown,
a clear hyperintense periventricular lesion
sits on the viewer's right of the image,
which by standard radiological convention
(``radiologist's view,'' patient-right on viewer-left)
corresponds to the patient's left brain.

\begin{figure}[htbp]
  \centering
  \includegraphics[width=0.30\textwidth]{vqarad_fc_102.jpg}
  \caption{VQA-RAD Case 2: axial brain MRI.
    The hyperintense periventricular lesion
    sits on the viewer's right of the image,
    i.e., the patient's \emph{left} hemisphere.}
  \label{fig:vqarad_case2}
\end{figure}

\begin{table}[htbp]
  \centering
  \caption{VQA-RAD Case 2: model responses.
    Expected answer: \emph{left}.}
  \label{tab:vqarad_case2_answers}
  \small
  \begin{tabular}{p{2.8cm}p{8cm}cc}
    \toprule
    Model & Response summary & Score & Type \\
    \midrule
    GPT-5.4-mini     & Locates lesion on the viewer's right
                       (visually correct), but misapplies the convention
                       $\to$ concludes patient's right                & 2 & faithfulness \\
    Gemini 2.5 Flash & Misreads the brain MRI as an abdominal CT;
                       describes gallstones in the gallbladder        & 0 & faithfulness \\
    Qwen3.5-35B      & Recites convention correctly, but \emph{mislocates}
                       the hyperintensity on the left of the image
                       $\to$ concludes patient's right                & 2 & faithfulness \\
    Qwen3-VL-235B    & Recites convention with extra detail
                       (DWI, ischemia), but mislocates the lesion
                       in ``right frontal/parietal'' $\to$ patient's right & 2 & faithfulness \\
    \bottomrule
  \end{tabular}
\end{table}

\textbf{All four models fail along three different paths.}
Gemini misidentifies the modality entirely,
describing an abdominal CT with gallstones on what is unambiguously
an axial brain MRI, which represents the score-0 extreme
of the bimodal distribution in \Cref{sec:vqarad}.
The other three converge on the same wrong answer (``patient's right'')
through two opposite error paths.
GPT-5.4-mini sees the lesion correctly on the viewer's right
but inverts the radiological convention,
equating ``right of the image'' with ``patient's right brain''
despite acknowledging the convention in passing.
The two Qwen models do the opposite:
they recite the convention correctly
(Qwen3-VL even prefaces its reasoning with
``the patient's right side appears on the left side of the image'')
but then mislocate the hyperintensity itself,
placing it on the image's left (Qwen3.5)
or directly in the patient's right hemisphere (Qwen3-VL).
A confident, clinically phrased ``right side''
thus masks two different underlying defects:
broken convention application in GPT,
broken visual localization in the Qwen pair.

\textbf{Linguistic prior overriding visual evidence.}
The pattern echoes POPE Case~2 (\Cref{sec:case2_pope})
at a higher level of abstraction.
In POPE, the co-occurrence prior
(\emph{outdoor + person + eating $\to$ backpack})
overrode visual evidence.
Here, the three non-Gemini models all default
to the most rehearsed clinical phrasing
(``the lesion is on the right side of the brain''),
each via a different local breakdown,
yet yielding the same fluent, confidently clinical wrong answer.
Qwen3-VL even posts
``\textbf{Final Answer: The lesion is on the RIGHT side of the brain}.''
The mechanism is structurally the same:
when visual evidence is perceptually demanding,
a strong language prior prevails.
The medical domain merely raises the stakes,
because the prior here is expert phrasing
that lends the wrong answer surface plausibility.

\subsubsection{Synthesis: Knowledge Without Perception}

Together the two cases delineate the two dominant failure modes
of general-purpose MLLMs in radiology:
under-claim in Case~1
(the panel describes the findings but refuses to convert them into a diagnosis,
producing a false negative on a screening question)
and over-claim in Case~2
(the panel produces the rehearsed clinical phrasing
even when the visual or convention step is wrong,
producing a false positive on the wrong hemisphere).
Both pathways are classified as faithfulness hallucinations by the judge,
which explains the near-monopoly of that category
in \Cref{tab:vqarad_results} ($>$96\% across all four models).

Both modes reflect a single deeper gap.
All four models can write fluently about
DWI sequences, periventricular white matter,
and the radiologist's-view convention,
linguistic content
that the language backbone has clearly absorbed.
What they cannot reliably do is decide
which half of a brain MRI contains a bright lesion,
a perception task that involves no medical knowledge at all.
When the visual signal is clear (Case~1),
conservative alignment suppresses the correct answer;
when it is weak (Case~2),
language priors fabricate one in its place.
Fluent medical knowledge, in short,
does not entail medical visual competence.
